\section{Determinism, idempotency, resumption}
\label{sec:determinism}

\spec{17.2} states three obligations that are usually treated as informal
quality properties. They are provable in this model, and two of them are
corollaries of \cref{sec:guarantees}.

\begin{theorem}[determinism]
\label{thm:determinism}
$\plan{\st}$, the versions $\vers{\st}{\cdot}$, the planned publish order and the ordered planning
diagnostic list are functions of $(\st, \Head, \mathrm{config})$ alone. Two
conforming implementations, on two machines, at two times, produce identical plan output. Concurrent task completion, logs, and other operational events are outside this claim.
\end{theorem}

\begin{proof}
Every quantity in \cref{sec:model,sec:propagation} is defined by a closed
expression in $\st$, \Head{} and the configuration:
\cref{def:window,def:sources} from ancestry and tags, \cref{def:delivery} and
the owed windows of \cref{def:unit} from the same;
\cref{def:unit} from commit messages and the configured types, separator and
limits; \cref{def:reach} from the graph at \Head{};
\cref{def:propchan,def:phase1,def:channels,def:resolvable,def:admission} from
those; \cref{def:plan,def:order} from those. None mentions wall-clock time,
commit dates, tag creation order, or the run history.

Three points admit a choice and each is resolved by a stated total order.
Conflicting propagated channels and disagreeing origin channels resolve to the
byte-wise least package name (\diag{W160}, \cref{def:propchan}). Competing
directives resolve by \spec{11.6} precedence: newest commit first, then last
unit within the commit, a total order on units, since commits carry a
topological index and units an index within their commit. Publish-order ties
resolve by byte-wise least name (\cref{prop:order}). \spec{17.2} additionally
requires every collection reaching the output to be sorted by package name,
then commit topological index, then unit index, so no unordered container's
iteration order is observable. A function of fixed inputs with all ties broken
by total orders is single-valued.
\end{proof}

\begin{theorem}[resumption after partial failure]
\label{thm:resumption}
Assume the retry invariant and $\wf{\st_1}$. Let $\st_2 = \Run{\st_1}{F}$
with blocked set $B = \blocked{\st_1}{F}$ and published set $\pubset{\st_1}{F}$,
and suppose the run reports \diag{E201} for no package. Then
\[
  \plan{\st_2} \;=\; F \cup B \cup A,
  \qquad
  A = \bigl\{\, p \in \pubset{\st_1}{F} \;:\; \owes{\st_2}{P}{p}
      \text{ for some } P \in F \cup B \,\bigr\},
\]
and $\vers{\st_2}{p} = \vers{\st_1}{p}$ for every $p \in F \cup B$.
\end{theorem}

\begin{proof}
$\subseteq$: $\plan{\st_2} \subseteq \plan{\st_1}$ by \cref{lem:plan-mono},
and every member of $\plan{\st_1}$ lies in $F$, in $B$ or in
$\pubset{\st_1}{F}$. Take $p \in \plan{\st_2} \cap \pubset{\st_1}{F}$. It is
tagged at \Head{} in $\st_2$, so by \cref{lem:tagged-target} its only cause is
a unit $u$ with $\admit{\st_2}{u}{p}$, and every $P \in \owed{\st_2}{u}{p}$ is
in $\plan{\st_2} \subseteq \plan{\st_1}$ and untagged at \Head{}, hence not
published in run 1, hence in $F \cup B$; so $p \in A$.
$\supseteq$: each $p \in F \cup B$ was planned in run 1 and not released, and
the run reported \diag{E201} for nothing, so \cref{thm:g3} with $k = 2$ gives
$p \in \plan{\st_2}$ at the same version. Each $p \in A$ has an admitted unit
with a non-empty owed set, so $\propb{\st_2}{p} \neq \bnone$, and $\held{p}$
is constant and false; so $p \in \plan{\st_2}$.
\end{proof}

The set $A$ is new with delivery: the consumers that proceeded past a commit
on a cause of their own while a provider failed or was blocked, and are owed
that provider's release. They return to the plan at a version their catch-up
computes, not at the version they were first planned at (\cref{thm:g3} says
nothing about a released package), and they are ordered after the provider
(\cref{lem:tagged-target}). Under the window test of CCME 2.0.0, $A$ was
empty, because a published package could not be planned again at the same
\Head{}; that is the orphan of \cref{ex:orphan}, second shape.

\Cref{thm:resumption} is the exact statement of what the operator is promised
after a failed run under the retry invariant: the next run does the rest of
the same job, at the same numbers for what never released and at the owed
catch-up for what proceeded, without coordinator state to clean up. An
ambiguous external publish outcome still requires \spec{19.4} reconciliation.
It is also what \spec{19.6} checks. After a run in which nothing failed or was
blocked, $F \cup B = \emptyset$, so $A = \emptyset$ and the re-planned set must
be empty; a non-empty one is \diag{E199}. Inside the retry invariant it
contradicts \cref{thm:resumption}; outside it, \spec{19.6} reports the changed
admission inputs for renewed review.

\begin{corollary}[idempotency]
\label{cor:idempotency}
Running the engine twice with no intervening commits and no publish yields the
same plan. Under the retry invariant, running it after a fully successful
publish from a well-formed state yields the empty plan.
\end{corollary}

\begin{proof}
The first is \cref{thm:determinism}: the state is unchanged, so the plan is.
The second is \cref{thm:resumption} with $F = \emptyset$: then $B = \emptyset$,
no \diag{E201} is reported (it needs a failed package), and $A = \emptyset$
because nothing is left to owe; so $\plan{\st_2} = \emptyset$. Equivalently,
every $p \in \plan{\st_1}$ is tagged at \Head{} in $\st_2$ together with every
package it was owed by, so \cref{lem:tagged-target} leaves it no cause; the
invariant excludes both input drift and newly visible catch-up.
\end{proof}

\begin{remark}[what determinism forbids]
\Cref{thm:determinism} is the reason \spec{13.11} may permit
optimisation of the computation but not of its result. An implementation may
cache identical traversals, invert loops, cache lexical parses and publish concurrently; it may
not introduce parallelism whose scheduling affects the iteration order of any
collection that reaches the output. \Cref{sec:complexity} records the precise
conditions under which such transformations preserve those observables.
\end{remark}
